\section[Categoría con clivaje definida por un pseudofuntor]{Categoría con clivaje definida por un pseudofuntor
$\mathcal{E}^\circ\to\Cat$} \label{VI.8}

Llamemos, para abreviar, \emph{pseudofuntor} de~$\cal{E}^\circ$
a~$\Cat$ (habría que decir pseudofuntor \emph{normalizado}) a un
conjunto de datos a), b), c) como los anteriores, que satisfacen las
condiciones $A')$ y~$B)$. En el número precedente hemos asociado a una
categoría con clivaje normalizado sobre~$\cal{E}$ un pseudofuntor
$\cal{E}^\circ\to\Cat$; aquí vamos a indicar la construcción inversa.
Dejaremos al lector la verificación de la mayor parte de los detalles, así
como del hecho de que estas construcciones son efectivamente inversas entre
sí. Con mayor precisión, convendría considerar los pseudofuntores
$\cal{E}^\circ\to\Cat$ como los objetos de una nueva categoría y demostrar
que nuestras construcciones proporcionan equivalencias, cuasi-inversas
entre sí, entre esta última y la categoría de las categorías con clivaje
sobre~$\cal{E}$ definida en el número precedente.

Pongamos
$$
\cal{F}_\circ=\underset{S\in\Ob(\cal{E})}\coprod
\ifthenelse{\boolean{orig}}
{\Ob\cal{F}(S),}
{\Ob(\cal{F}(S)),}
$$
la suma disjunta de los conjuntos
\ifthenelse{\boolean{orig}}
{$\Ob\cal{F}(S)$}
{$\Ob(\cal{F}(S))$}
(N.B. escribiremos aquí $\cal{F}(S)$ y no $\cal{F}_S$ para el valor en el
objeto $S$ de~$\cal{E}$ del pseudofuntor dado, a fin de evitar confusiones
de notación más adelante). Tenemos, pues, una aplicación evidente:
$$
p_\circ\colon\cal{F}_\circ\to\Ob\cal{E}.
$$
Sean
$$
\ifthenelse{\boolean{orig}}
{\xi}
{\overline{\xi}}
=(S,\xi),\quad \overline{\eta}=(T,\eta)\qquad\qquad
\text{(con $\xi\in\Ob\cal{F}(S)$, $\eta\in\Ob\cal{F}(T)$)}
$$
dos elementos de~$\cal{F}_\circ$, y sea $f\in\Hom(T,S)$; pondremos
$$
h_f(\overline{\eta},\overline{\xi})=\Hom_{\cal{F}(T)}(\eta,f^*(\xi)).
$$
\marginpar{176}%
Si además tenemos un morfismo $g\colon U\to T$ en~$\cal{E}$ y un
$\zeta\in\Ob\cal{F}(U)$, definimos una aplicación, denotada por
$(u,v)\mto u\circ v$:
$$
h_f(\overline{\eta},\overline{\xi})\times
h_g(\overline{\zeta},\overline{\eta})\to
h_{fg}(\overline{\zeta},\overline{\xi}),
$$
es decir, una aplicación
$$
\Hom_{\cal{F}(T)}(\eta,f^*(\xi))\times\Hom_{\cal{F}(U)}(\zeta,g^*(\eta))
\to \Hom_{\cal{F}(U)}(\zeta,(fg)^*(\xi)),
$$
mediante la fórmula
$$
u\circ v=c_{f,g}(\xi)\cdot g^*(u)\cdot v,
$$
es decir, $u\circ v$ es el compuesto de la sucesión
$$
\zeta
\ifthenelse{\boolean{orig}}
{\lto{u}}
{\lto{v}}
g^*(\eta)\lto{g^*(u)}
g^*f^*(\xi)\lto{c_{f,g}(\xi)} (fg)^*(\xi).
$$
Pondremos, por otra parte,
$$
h(\overline{\eta},\overline{\xi})=\underset{f\in\Hom(T,S)}\coprod
h_f(\overline{\eta},\overline{\xi}),
$$
y las aplicaciones de composición precedentes definen aplicaciones
$$
h(\overline{\eta},\overline{\xi})\times
h(\overline{\zeta},\overline{\eta})\to
h(\overline{\zeta},\overline{\xi}),
$$
mientras que la definición de los $h(\overline{\eta},\overline{\xi})$
implica una aplicación evidente:
$$
p_{\overline{\eta},\overline{\xi}}\colon
h(\overline{\eta},\overline{\xi})\to \Hom(T,S).
$$
Dicho esto, se verifican los puntos siguientes:
\begin{enumerate}
\item[1)] La composición entre elementos de los
\ifthenelse{\boolean{orig}}
{$h(\eta,\xi)$}
{$h(\overline{\eta},\overline{\xi})$}
es \emph{asociativa}.
\item[2)]
\marginpar{177}
Para todo
\ifthenelse{\boolean{orig}}
{$\overline{\xi}=(\xi,S)$}
{$\overline{\xi}=(S,\xi)$}
en~$\cal{F}_\circ$, consideremos el elemento
\ifthenelse{\boolean{orig}}
{}
{identidad}
de
$$
\ifthenelse{\boolean{orig}}
{h_{\id_S}(\overline{\xi},\overline{\xi})=
\Hom_{\cal{F}_S}(\id_S^*(\xi),\xi)=\Hom_{\cal{F}_S}(\xi,\xi),}
{h_{\id_S}(\overline{\xi},\overline{\xi})=
\Hom_{\cal{F}(S)}(\id_S^*(\xi),\xi)=\Hom_{\cal{F}(S)}(\xi,\xi),}
$$
y su imagen en $h(\overline{\xi},\overline{\xi})$. Este elemento es un
\emph{elemento neutro} por la izquierda y por la derecha para la
composición entre elementos de los
$h(\overline{\eta},\overline{\xi})$.

Esto muestra ya que \emph{se obtiene una categoría} $\cal{F}$ al poner
$$
\Ob\cal{F}=\cal{F}_\circ,\quad \Fl{\cal{F}}=
\ifthenelse{\boolean{orig}}
{\underset{\xi,\overline{\eta}\in\cal{F}_\circ} \coprod
h(\overline{\eta},\overline{\xi}).}
{\underset{\overline{\xi},\overline{\eta}\in\cal{F}_\circ}
\coprod h(\overline{\eta},\overline{\xi}).}
$$
(N.B. no se puede tomar simplemente como $\Fl\cal{F}$ la \emph{unión}
de los conjuntos $h(\overline{\eta},\overline{\xi})$, pues estos no son
necesariamente disjuntos). Además:
\item[3)] Las aplicaciones $p_\circ\colon \Ob\cal{F}\to\Ob\cal{E}$ y
$p_1=(p_{\overline{\eta},\overline{\xi}})\colon\Fl\cal{F}\to\Fl{\cal{E}}$
definen un \emph{funtor} $p\colon\cal{F}\to\cal{E}$. De este modo,
$\cal{F}$ se convierte en una categoría sobre~$\cal{E}$; además, la
aplicación evidente
$h_f(\overline{\eta},\overline{\xi})\to
\Hom(\overline{\eta},\overline{\xi})$ induce una \emph{biyección}
$$
h_f(\overline{\eta},\overline{\xi})\isomto
\Hom_f(\overline{\eta},\overline{\xi}).
$$
\item[4)] Las aplicaciones evidentes
$$ \Ob\cal{F}(S)\to\cal{F}_\circ=\Ob\cal{F},\qquad \Fl\cal{F}(S)\to
\Fl\cal{F},
$$
donde la segunda viene definida por las aplicaciones evidentes
$$
\Hom_{\cal{F}(S)}(\xi,\xi')=
h_{\id_S}(\overline{\xi},\overline{\xi}')\to
\Hom(\overline{\xi},\overline{\xi}')
$$
definen un \emph{isomorfismo}
$$
i_S \colon \cal{F}(S)\isomto\cal{F}_S.
$$
\item[5)] Para todo objeto $\overline{\xi}=(S,\xi)$ de~$\cal{F}$ y todo
morfismo $f\colon T\to S$ de~$\cal{E}$, consideremos
\marginpar{178}
el elemento $\overline{\eta}=(T,\eta)$ de~$\cal{F}_T$, con
$\eta=f^*(\xi)$, y el elemento $\alpha_f(\xi)$ de
$\Hom(\overline{\eta},\overline{\xi})$, imagen de~$\id_{f^*(\xi)}$ por
el morfismo
$\Hom_{\cal{F}(T)}(f^*(\xi),f^*(\xi))=
h_f(\overline{\eta},\overline{\xi})
\to \Hom_f(\overline{\eta},\overline{\xi})$. \emph{Este elemento es
cartesiano y es la identidad de $\overline{\xi}$ si} $f=\id_S$; en otros
términos, el conjunto de los $\alpha_f(\xi)$ define un
\emph{clivaje normalizado de~$\cal{F}$ sobre} $\cal{E}$. Además, por
construcción, es conmutativo el diagrama de funtores
$$
\xymatrix{
\cal{F}(S) \ar[r]^{f^*} \ar[d]_{i_S} & \cal{F}(T) \ar[d]^{i_T} \\
\cal{F}_S \ar[r]^{f^*_{\cal{F}}} & \cal{F}_T }
$$
donde $f^*_{\cal{F}}$ es el funtor imagen inversa por~$f$, relativo al
clivaje considerado sobre~$\cal{F}$. Finalmente:
\item[6)] las transformaciones naturales $c_{f,g}$ dadas con el
pseudofuntor son llevadas, mediante los isomorfismos $i_S$, a las
transformaciones naturales $c_{f,g}$ asociadas al clivaje de~$\cal{F}$.
\end{enumerate}

Nos limitamos a dar la verificación de 1) (que es, si acaso, menos trivial
que las demás). Basta probar la asociatividad de la composición entre los
\ifthenelse{\boolean{orig}}
{objetos}
{elementos}
de conjuntos de la forma $h_f(\overline{\eta},\overline{\xi})$.
Consideremos, pues, en~$\cal{E}$ morfismos
$$
S\lfrom{f} T\lfrom{g} U\lfrom{h} V
$$
y objetos
$$
\xi,\ \eta,\ \zeta,\ \tau
$$
en $\cal{F}(S),\cal{F}(T),\cal{F}(U),\cal{F}(V)$, respectivamente, y,
por último, elementos
$$
\begin{array}{ccc}
u\in
h_f(\overline{\eta},\overline{\xi})&=&\Hom_{\cal{F}(T)}(\eta,f^*(\xi))\\
v\in
h_g(\overline{\zeta},\overline{\eta})&=&\Hom_{\cal{F}(U)}(\zeta,g^*(\eta))\\
w\in
h_h(\overline{\tau},\overline{\zeta})&=&\Hom_{\cal{F}(V)}(\tau,h^*(\zeta)).
\end{array}
$$
Queremos
\marginpar{179}
probar la fórmula
$$
(u\circ v)\circ w = u\circ (v\circ w),
$$
que es una igualdad en
$\Hom_{\cal{F}(V)}(\tau,(fgh)^*(\xi))$. En virtud de las definiciones,
los dos miembros de esta igualdad se obtienen componiendo a lo largo de
los contornos superior e inferior, respectivamente, del diagrama
siguiente:
$$
\xymatrix{
\tau \ar[r]^w \ar[drr]_{v\circ w} & h^*(\zeta)
\ar[r]^{h^*(v)} \ar@/^1.4pc/[rrrrr]^{h^*(u\circ v)} & h^*g^*(\eta)
\ar[rr]^{h^*g^*(u)}
\ar[d]_{c_{g,h}(\eta)} & & h^*g^*f^*(\xi) \ar[rr]^{h^*(c_{f,g}(\xi))}
\ar[d]^{c_{g,h}(f^*(\xi))} & & h^*(fg)^*(\xi) \ar[d]^{c_{fg,h}(\xi)} \\
& & (gh)^*(\eta) \ar[rr]^{(gh)^*(u)} & & (gh)^*f^*(\xi)
\ar[rr]^{c_{f,gh}(\xi)} & & (fgh)^*(\xi) }
$$
Ahora bien, el cuadrado central es conmutativo porque $c_{g,h}$ es una
transformación natural, y el cuadrado de la derecha es conmutativo en
virtud de la condición~$B)$ para un pseudofuntor. De aquí se sigue el
resultado anunciado.

Naturalmente, queda por precisar, cuando el pseudofuntor considerado
proviene ya de una categoría con clivaje normalizado $\cal{F}'$
sobre~$\cal{E}$, cómo se obtiene un isomorfismo natural entre $\cal{F}'$
y~$\cal{F}$. Dejamos el detalle al lector.

Dejamos también al lector la interpretación, en términos de
pseudofuntores, de la noción de imagen inversa de una categoría con
clivaje $\cal{F}$ sobre~$\cal{E}$ por un funtor de cambio de base
$\cal{E}'\to\cal{E}$.
